\documentclass{article}
\usepackage[utf8]{inputenc}
\usepackage{geometry}
\geometry{a4paper, margin=1in}
\usepackage{minted}

\title{UAV Software Engineering: C++ Fundamentals}
\author{Topic 1.4: Modern C++ for Real-Time Robotics}
\date{}

\begin{document}

\maketitle

\section{Value Categories and Move Semantics}
In modern C++, understanding value categories (lvalues vs. rvalues) is crucial for performance. In robotics, transferring ownership of large objects---like a 640x480 \texttt{cv::Mat} frame or a dense point cloud---without copying the underlying memory buffer is vital for maintaining high frame rates.

\textbf{Move Semantics (\texttt{std::move}):} \texttt{std::move} casts an object to an rvalue reference (\texttt{T\&\&}), signaling that its resources can be "stolen" by the recipient.
\begin{itemize}
    \item \textbf{When it saves a copy:} Passing a large, pre-allocated \texttt{std::vector} into a class constructor. By using \texttt{std::move}, the internal pointer is transferred in $O(1)$ time, avoiding an $O(N)$ deep copy.
    \item \textbf{When it silently does nothing (or hurts):} Returning a local variable by value using \texttt{return std::move(local\_var);} is an anti-pattern. It disables the compiler's Return Value Optimization (RVO/NRVO), which would have eliminated the copy entirely by constructing the object directly in the caller's memory space.
\end{itemize}

\section{Modern Constructs (C++17 \& C++20)}
Real-time loops benefit immensely from modern C++ features that enforce compile-time checks and eliminate runtime overhead.

\subsection{Compile-time Evaluation (\texttt{constexpr})}
Instead of computing trigonometric lookup tables (e.g., sine/cosine for drone kinematics) during node initialization, \texttt{constexpr} allows you to compute and store these tables in the binary at compile time, eliminating initialization latency.

\subsection{Non-Owning Views (\texttt{std::span} and \texttt{std::string\_view})}
Passing raw pointers and sizes into functions is a classic source of out-of-bounds Undefined Behavior (UB). \texttt{std::span<T>} (C++20) acts as a lightweight, non-owning view over a contiguous sequence of memory (like a \texttt{std::vector} or a raw C-array). It provides bounds checking and iterator support without allocating memory or copying data.

\section{Lambdas and Capture Pitfalls}
ROS 2 heavily utilizes lambdas for asynchronous callbacks (e.g., processing a LiDAR scan). A critical failure mode in UAV software is the \textbf{dangling reference} inside a lambda.

\begin{itemize}
    \item \textbf{The Dangling \texttt{this} Pointer:} If a node captures \texttt{[this]} in a timer callback, but the node instance is destroyed before the timer finishes firing, dereferencing \texttt{this} causes a segmentation fault. (In ROS 2, managing object lifecycles with \texttt{std::shared\_ptr} mitigates this).
    \item \textbf{The Dangling Reference (\texttt{[\&]}):} Capturing local variables by reference and passing the lambda to a detached thread guarantees a crash. The thread will attempt to read memory that has already been popped off the stack.
\end{itemize}

\section{Time Management in Control Loops}
A PID controller calculates the derivative term as $de / dt$. If $dt$ is negative or artificially stretched, the quadrotor will become unstable and crash.
\begin{itemize}
    \item \textbf{\texttt{std::chrono::system\_clock}:} Represents the wall-clock time. It is subject to NTP (Network Time Protocol) synchronization and daylight saving time adjustments. \textbf{Never use this for control loops}, as it can jump backwards.
    \item \textbf{\texttt{std::chrono::steady\_clock}:} A monotonic clock that is guaranteed to never decrease. This is the only acceptable clock for measuring $dt$, loop sleep intervals (\texttt{std::this\_thread::sleep\_until}), and calculating velocities.
\end{itemize}

\section{Error Handling in Hot Paths}
\textbf{Exceptions (\texttt{throw/catch}) are strictly forbidden in real-time hot paths.} Exception unwinding is non-deterministic and can cause execution stalls lasting milliseconds.

Exceptions are only acceptable during initialization (e.g., failing to parse a configuration YAML or failing to open an I2C device node at startup). Inside the 400Hz flight loop, use modern alternative types:
\begin{itemize}
    \item \textbf{\texttt{std::optional<T>}:} Represents a value that may or may not exist (e.g., reading an I2C sensor where the bus might be busy, returning no data for one tick).
    \item \textbf{\texttt{std::variant<T, ErrorCode>}:} A type-safe union. Combined with \texttt{std::visit}, it forces the programmer to explicitly handle both the success and the failure state at compile time.
\end{itemize}

\section{STL Cost Model and Memory Arenas}
\subsection{Cache Locality and Contiguity}
The CPU fetches memory in cache lines (typically 64 bytes). 
\begin{itemize}
    \item \textbf{\texttt{std::vector}, \texttt{std::array}:} Elements are strictly contiguous. The CPU hardware prefetcher predicts the access pattern, resulting in near-zero cache misses.
    \item \textbf{\texttt{std::list}, \textbf{Maps/Sets}:} Node-based containers scatter data across the heap. Traversing them causes constant L1/L2 cache misses, resulting in massive latency spikes.
\end{itemize}

\subsection{Heap Allocation in Hot Loops}
Calling \texttt{new}, \texttt{malloc}, or \texttt{push\_back} (when it triggers a capacity reallocation) requires an OS lock to find free heap space. \textbf{Rule of thumb:} Pre-allocate everything during node initialization. If dynamic sizing is strictly required, use \texttt{std::pmr} (Polymorphic Memory Resources) to assign a pre-allocated stack arena to your containers, ensuring $O(1)$ allocation without kernel involvement.

\section{Atomics and the Memory Model}
To safely pass data between a 1000Hz sensor reading thread and a 50Hz ROS 2 publishing thread without mutexes (which cause blocking and jitter), we use lock-free data structures built on \texttt{std::atomic}.

\subsection{Basic Atomic Operations}
Interacting with atomic variables should always be done explicitly using the base methods rather than overloaded assignment operators:
\begin{itemize}
    \item \textbf{\texttt{.load(std::memory\_order)}:} Used to safely \textbf{read} the value of an atomic variable.
    \item \textbf{\texttt{.store(value, std::memory\_order)}:} Used to safely \textbf{write} a value to an atomic variable.
\end{itemize}
By explicitly passing a \texttt{memory\_order} parameter, you instruct the compiler and the CPU's instruction reorderer on exactly what level of synchronization is required, balancing safety and bare-metal performance.

\subsection{The Acquire-Release Pattern}
In a Single-Producer/Single-Consumer (SPSC) Ring Buffer, we must ensure that the producer's data is fully written to RAM before the consumer reads it. We enforce this through specific memory orders:
\begin{itemize}
    \item \textbf{\texttt{memory\_order\_relaxed}:} Offers no cross-thread synchronization. It simply guarantees that the read/write of this specific variable is atomic. \textit{When to use:} Use it when a thread is reading an atomic variable that it exclusively "owns" and modifies (e.g., the Producer reading its own \texttt{head\_} index, or the Consumer reading its own \texttt{tail\_} index).
    \item \textbf{\texttt{memory\_order\_acquire}:} Used exclusively with \texttt{.load()}. It acts as a one-way synchronization gate. \textit{When to use:} The Consumer uses it when reading the Producer's \texttt{head\_} (or the Producer reading \texttt{tail\_}). It guarantees that no memory operations (like reading the payload from the buffer) can be reordered by the CPU \textit{before} this \texttt{.load()}. It ensures the thread safely "acquires" all memory updates published by the other thread.
    \item \textbf{\texttt{memory\_order\_release}:} Used exclusively with \texttt{.store()}. It acts as the opposite one-way gate. \textit{When to use:} The Producer uses it when updating \texttt{head\_} to signal new data is ready (and the Consumer uses it when updating \texttt{tail\_} to free up space). It guarantees that all previous memory writes (i.e., writing the payload into the buffer array) are finished and visible in RAM \textit{before} the \texttt{.store()} executes. It safely "publishes" the data.
\end{itemize}

\section{Undefined Behavior (UB) and Sanitizers}
Undefined Behavior means the C++ standard imposes no requirements on the compiler. It may format your hard drive or silently corrupt telemetry. 

\textbf{Tooling:} Do not rely on code review alone to catch UB. Integrate these into your CMake/CI pipeline:
\begin{itemize}
    \item \textbf{ASan (AddressSanitizer):} Catches out-of-bounds array access, use-after-free, and memory leaks.
    \item \textbf{UBSan (UndefinedBehaviorSanitizer):} Catches signed integer overflow, division by zero, and unaligned pointers.
    \item \textbf{TSan (ThreadSanitizer):} Catches data races (reading and writing to the same memory concurrently without synchronization).
\end{itemize}

\textbf{SEI CERT C++ Rules (Example):}
Rule EXP54-CPP: \textit{"Do not access an object outside of its lifetime."} Returning a reference to a local variable or capturing a temporary by reference in an asynchronous lambda directly violates this standard.

\section{Benchmark: Lock-Free vs Mutex Queue}
To empirically validate the advantages of a lock-free Single-Producer/Single-Consumer (SPSC) ring buffer over a traditional \texttt{std::mutex}-based queue, a stress test of 250,000 elements was performed on an Intel Core i5-9600K. 

\subsection{Compilation and Reproduction}
The benchmark must be compiled with POSIX threads support and maximum compiler optimizations to prevent artificial overhead:
\begin{verbatim}
g++ -std=c++17 -Wall -Wextra -pthread -O3 bench_queue.cpp -o bench_perf
\end{verbatim}

\subsection{Peak Performance (-O3 Optimization)}
Compiled with the command above, the lock-free architecture demonstrates its superiority by eliminating OS context-switching and exploiting cache alignment.
\begin{itemize}
    \item \textbf{Lock-Free SPSC:} Throughput $\approx 29.68$ Million ops/sec | p99 Latency = $40.16\ \mu$s (40,162 ns)
    \item \textbf{Mutex Queue:} Throughput $\approx 4.33$ Million ops/sec | p99 Latency = $268.77\ \mu$s (268,775 ns)
\end{itemize}
The lock-free queue is nearly 7 times faster in throughput and almost 7 times more stable in worst-case latency compared to the mutex baseline.

\subsection{ThreadSanitizer (TSan) Overhead}
When compiled with \texttt{-fsanitize=thread}, the compiler injects heavy instrumentation to track memory accesses and detect data races. While throughput collapses due to the diagnostic overhead, the lock-free implementation remains mathematically safe (no data races reported) and still outperforms the mutex version:
\begin{itemize}
    \item \textbf{Lock-Free SPSC (TSan):} Throughput $\approx 721$ Thousand ops/sec | p99 Latency = $1.48$ ms (1,485,364 ns)
    \item \textbf{Mutex Queue (TSan):} Throughput $\approx 269$ Thousand ops/sec | p99 Latency = $5.90$ ms (5,900,952 ns)
\end{itemize}
This proves that the \texttt{std::memory\_order} constraints (acquire/release) were placed correctly, ensuring strict memory safety without sacrificing hard real-time performance.

\section{Timing Validation: OS Scheduler Jitter}
To simulate a real-time flight control loop, a fixed-rate execution test was developed using \texttt{std::chrono::steady\_clock} and \texttt{std::this\_thread::sleep\_until}. The benchmark measures the scheduling jitter (the deviation between the requested wake-up time and the actual wake-up time) introduced by the standard Ubuntu 22.04 kernel (CFS).

\subsection{Compilation}
\begin{verbatim}
g++ -std=c++17 -Wall -Wextra -pthread -O3 fixed_rate_loop.cpp -o fixed_rate_loop
\end{verbatim}

\subsection{Empirical Results (Idle CPU)}
\begin{itemize}
    \item \textbf{100 Hz Target (10 ms interval):} Avg Jitter $\approx 0.10$ ms | Max Overrun $\approx 0.20$ ms
    \item \textbf{400 Hz Target (2.5 ms interval):} Avg Jitter $\approx 0.09$ ms | Max Overrun $\approx 0.23$ ms
\end{itemize}

\subsection{Architectural Implications}
While an overrun of $\approx 230\ \mu$s is impressive for a general-purpose operating system, it is fundamentally inadequate for hard real-time robotics. At 400 Hz, a 0.23 ms scheduling delay consumes nearly 10\% of the available control window. Under heavy CPU load, generic OS schedulers can stall threads for multiple milliseconds, leading to catastrophic phase lag in PID controllers. This empirical evidence strictly mandates the adoption of a fully preemptive kernel (\texttt{PREEMPT\_RT}) for the flight controller deployment.
\section{Interview Questions}

\begin{enumerate}
    \item \textbf{Q: What is the difference between \texttt{std::unique\_ptr} and \texttt{std::shared\_ptr}?}\\
    \textbf{A:} \texttt{unique\_ptr} expresses exclusive ownership and has zero overhead compared to a raw pointer. \texttt{shared\_ptr} expresses shared ownership and uses a heap-allocated control block for atomic reference counting, which introduces memory and performance overhead.

    \item \textbf{Q: Why must a ROS 2 timer callback avoid executing long-blocking operations?}\\
    \textbf{A:} If using a single-threaded executor, a blocking callback prevents all other callbacks in the queue (including critical telemetry and safety checks) from executing, leading to node starvation and potential failsafe triggers.

    \item \textbf{Q: Explain why we use \texttt{std::chrono::steady\_clock} for PID timing.}\\
    \textbf{A:} \texttt{steady\_clock} is monotonic and immune to system time adjustments (like NTP syncs). If \texttt{system\_clock} is adjusted backwards, $dt$ becomes negative, mathematically destabilizing the PID derivatives.

    \item \textbf{Q: What is Perfect Forwarding?}\\
    \textbf{A:} It is a template pattern using \texttt{T\&\&} (forwarding reference) and \texttt{std::forward<T>} to pass arguments to another function while perfectly preserving their value category (lvalue or rvalue) and \texttt{const} qualifiers.

    \item \textbf{Q: You are writing a lock-free Single-Producer/Single-Consumer queue. Which memory ordering is required for the \texttt{push} and \texttt{pop} indices?}\\
    \textbf{A:} The producer must use \texttt{memory\_order\_release} when updating the write index to ensure the data is visible. The consumer must use \texttt{memory\_order\_acquire} when reading the write index to safely access the newly written data without race conditions.
\end{enumerate}



\end{document}